\chapter{Methodology}\label{chap:Methodology}

\section{Chapter Overview}\label{sec:method_overview}

This chapter explains how the proposed system is built. The system has two main parts. The first part is the emotion recognition module, which reads short webcam clips of the learner and predicts academic affective states. The second part is the adaptation module, which uses these predictions to choose the next teaching action and the next question. The two parts are connected by a small interface that turns the emotion module output into a stable input for the adaptation policy.

The present chapter focuses on the first part, the Facial Expression Recognition (FER) module. The adaptation module is described later in its own section, because it has different design concerns and different evaluation metrics. Section~\ref{sec:method_dataset} introduces the data and the affect labels. Section~\ref{sec:method_preparation} explains the data preparation. Section~\ref{sec:method_architecture} describes the proposed FER architecture. Sections~\ref{sec:method_training} and \ref{sec:method_metrics} then cover training and the evaluation protocol. Sections~\ref{sec:method_design_rationale}, \ref{sec:method_reproducibility}, and \ref{sec:method_integration} cover the design rationale, list the implementation details, and explain how the module connects to adaptive learning. The quantitative results of this FER module on DAiSEE are reported in Chapter~\ref{chap:results}.

The proposed FER architecture is shown in Figure~\ref{fig:fer_pipeline}. A short video clip is first sampled and preprocessed. It is then passed through a deep visual feature extractor, a channel attention block, a temporal model, and a temporal attention layer. The output is one independent score per affective label. The reasoning behind this exact pipeline is explained step by step in the following sections.

\begin{figure}[H]
    \centering
    \includegraphics[width=\textwidth]{figures/fer_pipeline.png}
    \caption{End-to-end FER architecture used for multilabel emotion recognition. The pipeline starts from a learner video clip, applies frame sampling and preprocessing, extracts spatial features with ResNet50, recalibrates them with a Squeeze-and-Excitation block, encodes their temporal pattern with a two-layer BiLSTM, pools time steps with temporal attention, and produces four independent sigmoid scores for boredom, engagement, confusion, and frustration. Per-label thresholds turn these scores into final binary affect labels.}
    \label{fig:fer_pipeline}
\end{figure}

\section{Dataset and Affect Labels}\label{sec:method_dataset}

The FER module is trained and evaluated on the DAiSEE dataset \cite{gupta_daisee_2016}. DAiSEE is a public dataset of about $9{,}068$ short webcam clips collected from $112$ students in study-like conditions, where each clip is annotated with four affective states: engagement, boredom, confusion, and frustration \cite{gupta_daisee_2016}. The original annotation scheme assigns each state an ordinal level from $0$ to $3$, but several recent FER works on DAiSEE simplify these levels into binary present/absent decisions because the most useful pedagogical information is whether a state is active during the clip rather than its exact intensity \cite{rao_recognition_2021,leong_deep_2020,abedi_improving_2021}.

Engagement, boredom, confusion, and frustration are kept as the target labels because they are the academic emotions that have repeatedly been shown to matter most during learning \cite{dmello_language_2012,pekrun_control-value_2006,tyng_influences_2017}. They are not the same as the basic emotions of happiness or fear that classical FER usually targets, and educational FER systematic reviews highlight that detecting them, rather than basic emotions, is what makes a model useful for tutoring and online classrooms \cite{lek_academic_2023,florestiyanto_emotion_2024,vistorte_integrating_2024}. DAiSEE was therefore selected as the main benchmark in this thesis because it is the standard public benchmark for these four states \cite{lek_academic_2023}.

A multilabel formulation is used, meaning that a single clip can be marked positive for several states at the same time, instead of being assigned to exactly one class. This matches what is observed in real DAiSEE annotations, where a learner can look both engaged and slightly confused, or both bored and slightly frustrated, in the same short clip \cite{gupta_daisee_2016}. Educational FER reviews also note that academic affect tends to overlap rather than appear in isolated single classes, which makes single-label classification a poor fit for this task \cite{lek_academic_2023,florestiyanto_emotion_2024}.

The dataset is split into the official DAiSEE training, validation, and test sets, exactly as released by Gupta et al.\ \cite{gupta_daisee_2016}. No clips are removed and no new split is created: the training set contains $5{,}358$ clips, the validation set contains $1{,}429$ clips, and the test set contains $1{,}784$ clips. The training set is used to fit the model parameters, the validation set is used to choose the best epoch and to tune per-label decision thresholds, and the test set is held out and only used at the end to report final metrics. This three-way split is the standard protocol used by recent DAiSEE studies \cite{abedi_improving_2021,mehta_three-dimensional_2022,su_leveraging_2024}.

\section{Data Preparation}\label{sec:method_preparation}

Each DAiSEE clip is first decoded into image frames in an offline preprocessing step. The decoder samples each ten-second clip at a low frame rate, producing a fixed sequence of $16$ frames per clip, which are stored as JPEG images on disk. The downstream FER model, however, expects a temporal axis of length $T = 24$, so a second per-batch sampling step is applied at training, validation, and test time. For each clip, $24$ indices are produced with \texttt{numpy.linspace} between $0$ and $15$, then cast to integers, which yields a deterministic sequence covering all $16$ source frames with a small number of repeated indices. Because the source clip already has fewer frames than the target length, the random starting-window logic for training-time temporal jitter is never actually triggered in the present setup; sampling is therefore the same at training and inference. The frames are loaded at a resolution of $224 \times 224$ pixels.

The choice of $T = 24$ was made for two reasons. First, academic affective states such as boredom or confusion build up over several seconds, so a single frame is not enough \cite{abedi_improving_2021,asaju_adaptive_2022}. Second, very long sequences make training slow and unstable on commodity GPUs, especially when a deep backbone is used per frame \cite{asaju_affects_2021}. A fixed length of $24$ time steps gives the BiLSTM a stable input shape across batches while keeping memory and time manageable.

Each frame is decoded with OpenCV, converted from BGR to RGB, resized to $224 \times 224$, normalised by dividing pixel values by $255$, and then standardised using the ImageNet mean $(0.485, 0.456, 0.406)$ and standard deviation $(0.229, 0.224, 0.225)$. ImageNet normalisation is used so that the inputs match the statistics of the data on which the ResNet50 backbone was originally pre-trained \cite{he_resnet_2016}. This is the standard preprocessing recipe followed by educational FER models that also start from ImageNet-pretrained backbones \cite{asaju_adaptive_2022,gupta_edfa_2023,aly_revolutionizing_2025}.

A small set of frame-level data augmentations is applied only at training time, never at validation or test time. Each frame can be horizontally flipped with probability $0.5$. Brightness and contrast can be jittered with probability $0.5$, by sampling a contrast factor uniformly between $0.85$ and $1.15$ and a brightness shift uniformly between $-15$ and $15$ pixel values. A small random rotation between $-7^{\circ}$ and $+7^{\circ}$ is applied with probability $0.25$, with reflective borders. Mild Gaussian blur with a $3 \times 3$ kernel is applied with probability $0.15$. Finally, a small random crop is applied with probability $0.25$, by sampling a crop ratio uniformly between $0.90$ and $1.0$ of the frame size and resizing back to $224 \times 224$. These augmentations are deliberately conservative. They simulate the kind of variation a webcam may introduce in a real classroom, such as small head rotations, lighting changes, slight motion blur, and partial face crops, which are exactly the situations that recent educational FER reviews report as the main practical sources of error \cite{lek_academic_2023,florestiyanto_emotion_2024,vistorte_integrating_2024}. Aggressive augmentations such as large rotations or strong colour distortions were avoided because they tend to destroy the very subtle facial cues that distinguish confusion and frustration.

A controlled oversampling step is also applied to the training set before training. The DAiSEE training labels are heavily imbalanced: most clips are labelled engaged, while clips marked confused or frustrated are far less common, a problem that almost every recent DAiSEE study has reported and tried to address \cite{gupta_daisee_2016,abedi_improving_2021,mehta_three-dimensional_2022,santoni_convolutional_2023,rianto_engagement_2025}. Oversampling means duplicating samples from minority groups so that the model sees them more often during training, while undersampling means dropping samples from majority groups; the literature on imbalanced learning has shown that oversampling tends to hurt minority recall less than undersampling in this kind of small-data setting \cite{santoni_convolutional_2023}.

The oversampling is applied on a dominant-group basis. Each training clip is assigned a single dominant group based on the priority order frustration, confusion, boredom, engagement; clips with no positive label fall into a separate group. Clips in the engagement group are kept once, clips in the boredom group are duplicated twice, clips in the confusion group three times, and clips in the frustration group four times. The duplication factors grow with how rare the group is, so that the rarest dominant states are seen most often during training. The full duplicated set is then shuffled with a fixed random seed for reproducibility. This is a controlled form of oversampling because the multipliers are fixed in advance and chosen from the original DAiSEE statistics, rather than letting an automatic resampler decide them. It is also less disruptive than synthetic methods like SMOTE that interpolate between samples in feature space \cite{santoni_convolutional_2023}, since here the original frames are reused without modification, and any acquisition variability is added back through the augmentation pipeline above. Figure~\ref{fig:fer_oversampled_dist} shows the resulting positive-clip count for each label after this controlled oversampling has been applied to the DAiSEE training set.

\begin{figure}[H]
    \centering
    \includegraphics[width=0.78\textwidth]{figures/fer_oversampled_label_dist.png}
    \caption{Number of positive training clips per affective label in the DAiSEE training set after the controlled oversampling described in Section~\ref{sec:method_preparation}. Engagement remains the most frequent label even after oversampling, but the rare states (confusion and frustration) now have much higher representation than in the raw DAiSEE training set.}
    \label{fig:fer_oversampled_dist}
\end{figure}

\section{Proposed FER Architecture}\label{sec:method_architecture}

The proposed architecture, summarised in Figure~\ref{fig:fer_pipeline}, takes a tensor of shape $[B, T, C, H, W]$ as input, where $B$ is the batch size, $T = 24$ is the length of the temporal axis fed to the model (built from the $16$ decoded frames per clip with the linear-interpolation sampling described in Section~\ref{sec:method_preparation}), $C = 3$ is the number of colour channels, and $H = W = 224$ is the spatial resolution. The output is a vector of four sigmoid scores per clip, one for each of boredom, engagement, confusion, and frustration. The pipeline contains five blocks: a spatial backbone, a channel attention block, a temporal recurrent encoder, a temporal attention pooler, and a multilabel classification head. Each block is described below, in the same order as in Figure~\ref{fig:fer_pipeline}.

\subsection{Spatial Backbone: ResNet50}\label{sec:method_resnet}

The spatial backbone is ResNet50 \cite{he_resnet_2016}. ResNet50 is a Convolutional Neural Network with $50$ layers that uses residual connections, which are shortcut paths that skip over groups of layers and add the input of a block back to its output \cite{he_resnet_2016}. Without these connections, very deep networks become hard to train, because the gradients used to update the early layers shrink as they pass through many layers \cite{he_resnet_2016}. The residual design fixes this and allows much deeper networks to be trained reliably, which is why ResNet50 has become a standard image feature extractor for many vision tasks, including educational FER \cite{asaju_adaptive_2022,asaju_affects_2021,gupta_edfa_2023,aly_revolutionizing_2025}.

ResNet50 is loaded with weights that were pre-trained on ImageNet and is fine-tuned end-to-end during training, rather than frozen. Pre-training means that the network has already learned generic visual filters, such as edges, textures, parts of faces, and lighting patterns, before this project starts \cite{he_resnet_2016}. Fine-tuning means that these filters are then adjusted on DAiSEE so that they capture cues that matter for academic affect, such as eyebrow shape, eye opening, mouth tension, or head orientation. This transfer-learning strategy is standard in educational FER, where datasets are too small to train such a deep network from random initialisation \cite{komaravalli_detecting_2023,gupta_edfa_2023,aly_revolutionizing_2025}.

To use ResNet50 as a pure feature extractor, its original $1000$-class ImageNet head is removed, leaving the convolutional body and the global pooling layer. For each input frame, the body produces a $2048$-dimensional feature vector that summarises the appearance of that frame. The $T = 24$ frames of a clip are passed through the backbone in a single batched call by reshaping the input to $[B \cdot T, C, H, W]$, and the resulting features are reshaped back to $[B, T, 2048]$ before the next block. This design keeps GPU usage efficient and matches how previous CNN-plus-temporal models on DAiSEE process clips frame by frame \cite{abedi_improving_2021,asaju_adaptive_2022,shiri_detection_2024}.

\subsection{Squeeze-and-Excitation Channel Attention}\label{sec:method_se}

After ResNet50, the per-frame feature vectors pass through a Squeeze-and-Excitation (SE) block, adapted from the original Squeeze-and-Excitation Networks design \cite{hu_squeeze_2018}. An SE block is a small attention module that learns, for each frame, which of the $2048$ feature channels are more useful and which can be down-weighted \cite{hu_squeeze_2018}. It does this with two small fully connected layers: the first reduces the $2048$-dimensional vector to a hidden size of $\max(2048 / 16, 64) = 128$ with a ReLU activation, and the second projects it back to $2048$ with a sigmoid activation, producing one importance weight per channel. These weights are then multiplied element-wise back onto the original features, so each channel is rescaled by its learned importance.

The reason for using SE attention here is that ResNet50 was trained on ImageNet, where the most informative channels are not necessarily those that matter for facial affect in a learning context \cite{he_resnet_2016}. The SE block lets the network re-weight ResNet50 features so that channels that respond to confusion-, boredom- or frustration-relevant cues are emphasised, and noisy ones are suppressed. Recent educational FER systems have used closely related channel-and-spatial attention modules, such as CBAM, on top of similar backbones for the same reason \cite{aly_revolutionizing_2025,liu_dynamic_nodate}. Compared with full CBAM, the SE variant used here is lighter, since it acts only on channels and does not add a spatial attention sub-module, which keeps the model small enough to be trained on a single GPU.

\subsection{Temporal Encoder: Two-Layer BiLSTM}\label{sec:method_bilstm}

The frame-level features are then passed to a two-layer Bidirectional Long Short-Term Memory network, or BiLSTM. A Long Short-Term Memory network is a type of Recurrent Neural Network designed to keep useful information from earlier time steps and forget less useful information through learned gates \cite{hochreiter_lstm_1997}. A bidirectional version reads the same sequence both forward and backward and concatenates the two hidden states at each time step, so each frame is encoded with knowledge of both the past and the future of the clip \cite{hochreiter_lstm_1997,asaju_adaptive_2022,tulsani_realtime_2025}. This is useful for academic affect, because cues like a delayed nod after a confusing question, or a short smile after a difficult problem, only become meaningful when the frames before and after are also taken into account.

The BiLSTM in this thesis has two stacked layers with hidden size $256$ in each direction, dropout $0.5$ between stacked layers, and reads the $[B, T, 2048]$ sequence produced by the SE block. Its output has shape $[B, T, 512]$, where $512 = 2 \cdot 256$ comes from concatenating the forward and backward hidden states. Stacking two layers, instead of one, lets the network represent affective patterns at different time scales: the lower layer can capture short, frame-to-frame transitions, while the upper layer can capture slower trends across the whole clip. Recent DAiSEE pipelines that combine a CNN backbone with a BiLSTM have shown that this kind of temporal recurrence is one of the strongest single design choices for engagement and affect recognition \cite{asaju_adaptive_2022,asaju_affects_2021,shiri_detection_2024,tulsani_realtime_2025}.

\subsection{Temporal Attention Pooling}\label{sec:method_temporal_attention}

A clip-level prediction must be made from $T = 24$ time-step vectors. A simple way would be to take the last hidden state of the BiLSTM, but this gives all the influence to the last frames of the clip and ignores the rest. Another simple option would be to average all $T$ vectors, but this gives equal weight to every frame, including ones where the face is occluded, the learner blinks, or the head turns away from the screen. A temporal attention layer is used instead.

The temporal attention layer is a small two-step network adapted from the general attention mechanism used in modern sequence models \cite{vaswani_attention_2017,mehta_three-dimensional_2022,su_leveraging_2024}. For each of the $T$ time-step vectors, it computes an unnormalised score with a fully connected layer of hidden size $128$ and a $\tanh$ activation, followed by another fully connected layer that produces a single scalar. The $T$ scalars are then turned into non-negative weights that sum to one with a softmax across time, and the final clip representation is the weighted sum of the time-step vectors with these weights. The output therefore has shape $[B, 512]$ and represents a weighted summary of the whole clip, where each frame contributes according to how relevant it is for the affective decision. The same idea, often called attention pooling, has been used in other DAiSEE-based engagement detectors and tends to focus on the most expressive frames in the clip \cite{mehta_three-dimensional_2022,su_leveraging_2024,liu_dynamic_nodate}.

\subsection{Multilabel Classification Head}\label{sec:method_head}

The clip representation is passed through a small classification head that outputs one logit per affective label. The head is a sequence of dropout, a linear projection to $256$ units, a ReLU activation, another dropout, and a final linear projection to $4$ logits. Both dropout layers use a dropout rate of $0.5$, which randomly zeroes half of the units during training and is a standard regulariser used in many recent DAiSEE pipelines \cite{asaju_adaptive_2022,shiri_detection_2024}.

A sigmoid activation is then applied to each of the four logits independently, producing a value between $0$ and $1$ for each label. These values are interpreted as the probability that the corresponding affective state is active in the clip. Because the sigmoid is applied to each label separately, the four predictions do not need to sum to one, which is exactly what multilabel classification requires \cite{rianto_engagement_2025}. This is different from a softmax head, which would force the four labels to compete with each other for the same probability mass and would make it impossible to predict, for example, that a learner is both engaged and slightly confused at the same time. The four sigmoid scores correspond, in order, to boredom, engagement, confusion and frustration, matching the DAiSEE label set \cite{gupta_daisee_2016}. They are turned into final binary predictions in the threshold-tuning step described in Section~\ref{sec:method_training}.

\section{Training Strategy}\label{sec:method_training}

The model is trained end-to-end. Both the ResNet50 backbone and the upper temporal blocks are updated by gradient descent. End-to-end training means that the visual feature extractor is allowed to specialise on educational facial cues during training, rather than staying fixed on generic ImageNet features. Recent educational FER work that fine-tunes the backbone, instead of freezing it, has reported clear gains compared with frozen-backbone variants on DAiSEE-style data \cite{abedi_improving_2021,gupta_edfa_2023,aly_revolutionizing_2025}.

\subsection{Loss Function and Class Imbalance}\label{sec:method_loss}

The loss function is the per-label binary cross-entropy with logits, with a positive-class weight per label. Binary cross-entropy is the standard loss for binary decisions, and applying it independently to each of the four labels turns it into a natural multilabel loss \cite{rianto_engagement_2025}. The per-label positive-class weight is computed from the training set as the ratio of negatives to positives for that label, then clipped to the interval $[0.5, 6.0]$ so that a single very rare class cannot dominate the gradient. This mirrors the weighted cross-entropy strategies used by other DAiSEE works that explicitly target class imbalance \cite{abedi_improving_2021,mehta_three-dimensional_2022,rianto_engagement_2025}.

This weighted loss complements the controlled oversampling described in Section~\ref{sec:method_preparation}: oversampling makes the model see minority groups more often, while the per-label weight makes a missed positive on a minority label more costly than a missed negative. The two together follow the imbalance-aware design pattern recommended by recent DAiSEE studies and by educational FER reviews more broadly \cite{abedi_improving_2021,mehta_three-dimensional_2022,santoni_convolutional_2023,rianto_engagement_2025,lek_academic_2023}.

\subsection{Optimisation}\label{sec:method_optim}

The model is optimised with AdamW. AdamW is a popular adaptive optimiser that keeps a running estimate of the gradient and its variance for each parameter and applies a separate weight-decay term. It is also the optimiser used by many recent DAiSEE pipelines \cite{shiri_detection_2024,tulsani_realtime_2025}. The learning rate is set to $5 \times 10^{-5}$, which is a small value chosen so that the pre-trained backbone is updated gently and does not lose its ImageNet knowledge. The weight-decay coefficient is $10^{-5}$. The batch size is $4$, which is a small value imposed by the high memory cost of running a $50$-layer backbone on $24$ frames per clip. Gradients are clipped at a maximum global norm of $5.0$ to avoid occasional large updates that could destabilise training, a common safeguard in recurrent architectures.

A learning-rate scheduler of type \texttt{ReduceLROnPlateau} watches the validation Macro F1 score. If the Macro F1 does not improve for two consecutive validation rounds, the learning rate is multiplied by $0.5$. The model is trained for up to $12$ epochs, with an early-stopping rule that interrupts training if the validation Macro F1 does not improve for four epochs in a row. The best checkpoint is the one with the highest validation Macro F1, not the one from the last epoch. This validation-driven model selection is the same protocol followed by recent DAiSEE engagement detectors that report careful imbalance-aware evaluation \cite{abedi_improving_2021,shiri_detection_2024,su_leveraging_2024}.

\subsection{Per-Label Threshold Tuning}\label{sec:method_thresholds}

Once the sigmoid scores have been computed, a threshold must be chosen for each label to produce a binary decision. A common default is to use $0.5$ for every label, but this default is rarely optimal in imbalanced multilabel settings, because it does not account for the different prior probabilities of different labels \cite{rianto_engagement_2025}. A per-label threshold tuning step is therefore included in the training pipeline.

After every epoch, the validation set is scored once with the current model, producing a probability matrix of size $N_{\mathrm{val}} \times 4$. For each of the four labels, a grid of candidate thresholds is searched over the values $0.10, 0.15, 0.20, \dots, 0.70$, in steps of $0.05$. For each candidate, the corresponding binary predictions are formed and the per-label F1 score is computed against the ground truth. The threshold that maximises the per-label F1 is kept for that label. The four chosen thresholds are then reused on the validation set to compute the global metrics that are reported each epoch. The thresholds saved with the best checkpoint are the ones used at test time. This procedure is consistent with how prior DAiSEE works adapt their decision thresholds when class imbalance is present, rather than fixing them at $0.5$ \cite{rianto_engagement_2025,mehta_three-dimensional_2022}.

\section{Evaluation Metrics}\label{sec:method_metrics}

Reporting only one global accuracy is not enough on this task, because accuracy can be misleading on imbalanced multilabel data, where a model that always predicts the majority labels can already look competitive \cite{lek_academic_2023,rianto_engagement_2025}. The evaluation in this thesis therefore uses several complementary metrics, all of them computed with the per-label thresholds described in Section~\ref{sec:method_thresholds}.

The first set is the F1 family. The F1 score is the harmonic mean of precision and recall: precision is the fraction of predicted positives that are actually positive, and recall is the fraction of real positives that are correctly retrieved \cite{rianto_engagement_2025}. Macro F1 is the average of the per-label F1 scores, where each label contributes equally regardless of its frequency. Weighted F1 is the same average but each per-label F1 is weighted by the number of real positives for that label, so frequent labels count more. Micro F1 is computed by pooling all the predictions across labels and computing one global F1 from those pooled counts. Samples F1 is computed per clip, by treating each clip as its own multilabel example and averaging the resulting F1 over clips. These four variants are the same family of multilabel F1 metrics used by recent DAiSEE-style and multilabel FER works \cite{rianto_engagement_2025,fu_integrating_2025,mehta_three-dimensional_2022}.

A complementary metric is exact-match accuracy, which counts a prediction as correct only when all four labels are correctly predicted for the clip. This is a strict criterion: a clip with three labels right and one label wrong is counted as a full miss. It is reported because it shows how often the system is fully aligned with the ground truth, which is a relevant quality for downstream adaptation that consumes the full label vector. Per-label accuracy is also reported for each of the four labels, since it gives a class-by-class view of how the model behaves and helps diagnose which labels are still hard.

\section{Design Rationale and Justification}\label{sec:method_design_rationale}

This section explains the reasoning behind the main design choices in the proposed FER pipeline. The justifications presented here are conceptual rather than experimental: they are grounded in what the educational FER literature reports as effective designs, in the structural properties of the DAiSEE dataset, and in the practical engineering constraints of training on a single consumer-grade GPU. The empirical evidence supporting these choices, including a side-by-side comparison of the selected configuration against alternative backbones, alternative temporal modules, and ablations of the imbalance-handling components, is reported separately as part of the architecture-selection experiments in Section~\ref{sec:results_fer_arch_selection}.

\textbf{Why ResNet50 rather than a deeper or transformer-based backbone.} ResNet50 is used because it provides a strong visual feature extractor at a controlled compute and memory cost \cite{he_resnet_2016}. Larger CNN backbones such as ResNet-152 or EfficientNetV2-L have been used on DAiSEE \cite{asaju_adaptive_2022,shiri_detection_2024}, and Vision Transformer or 3D-attention models have also been reported to improve four-class engagement detection \cite{mehta_three-dimensional_2022,su_leveraging_2024,liu_dynamic_nodate}. These models typically require more memory, more careful training schedules, or more training data than DAiSEE provides, and recent reviews observe that their gain on DAiSEE is small relative to the additional deployment cost \cite{lek_academic_2023,vistorte_integrating_2024}. ResNet50 with channel attention therefore sits in the middle of this accuracy--cost trade-off and is a conceptually sound default for an educational FER pipeline that should also run on standard hardware.

\textbf{Why a BiLSTM rather than a 3D CNN or temporal transformer.} A BiLSTM is used because it is an explicit sequence model that can be combined with a strong 2D backbone without doubling the parameter count, and because BiLSTM-based DAiSEE pipelines remain among the most reliable temporal designs in the literature \cite{asaju_adaptive_2022,asaju_affects_2021,tulsani_realtime_2025,shiri_detection_2024}. 3D CNNs such as C3D or 3D DenseNet, and temporal transformers, have been proposed for DAiSEE \cite{gupta_daisee_2016,mehta_three-dimensional_2022,su_leveraging_2024,liu_dynamic_nodate}, but they require more memory per clip and are harder to train stably at the small batch sizes that fit on a single consumer-grade GPU.

\textbf{Why multilabel rather than single-label classification.} A single-label formulation would force the model to choose only one of the four affective states per clip, which contradicts both the DAiSEE annotations and the way academic affect tends to co-occur in real clips \cite{gupta_daisee_2016,dmello_language_2012}. Recent reviews of educational FER also recommend treating affect as several co-occurring states rather than as a single mutually exclusive class \cite{lek_academic_2023,florestiyanto_emotion_2024}. The multilabel design adopted here is therefore the formulation that matches both the underlying labels and the pedagogical reality.

\textbf{Why imbalance-aware training with both oversampling and weighted loss.} Imbalance is not a side issue in DAiSEE; it is a defining property of the dataset \cite{gupta_daisee_2016,abedi_improving_2021}. Several DAiSEE-focused studies have shown that, without explicit imbalance handling, models collapse onto the majority class and produce nearly unusable minority-class metrics \cite{santoni_convolutional_2023,rianto_engagement_2025}. The combined design used here---controlled oversampling at the data level, weighted binary cross-entropy at the optimisation level, and per-label threshold tuning at the decision level---addresses imbalance at three different stages, which is more robust than relying on any single mechanism.

\textbf{Why per-label thresholds rather than the default $0.5$.} The four labels have very different prior probabilities in DAiSEE. A fixed threshold of $0.5$ for all labels would suppress too many positives on rare labels and would collapse minority recall. Per-label thresholds, tuned on the validation set, are a simple and inexpensive way to keep the model usable on rare labels without retraining, in line with the imbalance-aware decision strategies used by recent DAiSEE pipelines \cite{rianto_engagement_2025,mehta_three-dimensional_2022}.

\textbf{Webcam-only design.} The whole pipeline takes only a video clip as input. There are no physiological sensors, no microphone, and no extra wearable. This is a deliberate constraint. Recent reviews of educational FER and emotion-aware tutoring have repeatedly noted that multimodal sensor stacks are expensive, fragile, and hard to deploy at scale in real classrooms or online courses \cite{vistorte_integrating_2024,fernandez-herrero_evaluating_2024,florestiyanto_emotion_2024}. A webcam-only system runs on the hardware that online learners already have, which is the main reason for keeping the FER module unimodal in this thesis.

These design choices were also validated empirically through a series of architecture-selection experiments in which alternative backbones, alternative temporal modules, and ablations of the imbalance-handling components were trained under matched protocols and compared against the configuration described above; the corresponding measurements are reported in Section~\ref{sec:results_fer_arch_selection}.

\section{Implementation and Reproducibility Details}\label{sec:method_reproducibility}

All experiments are implemented in Python with PyTorch and torchvision. Frame I/O uses OpenCV, and metrics are computed with scikit-learn. The dataset class, model class, training script, and evaluation script are organised in separate files inside the experiment folder and are kept self-contained, so that the same training run can be reproduced from a single command. The frame extraction step and the construction of the oversampled training CSV are run once offline before training begins, and their outputs are stored as image folders and CSV files. This separation ensures that frame sampling, augmentation, and model fitting can be re-run without redoing the heavier video decoding step. Training and evaluation were run on a single NVIDIA Tesla T4 GPU on the Kaggle Notebooks platform.

The full set of hyperparameters used in this work is summarised in Table~\ref{tab:fer_hparams}.

\begin{table}[H]
\centering
\caption{FER training hyperparameters.}
\label{tab:fer_hparams}
\begin{tabular}{ll}
\toprule
\textbf{Hyperparameter} & \textbf{Value} \\
\midrule
Image size                      & $224 \times 224$ \\
Frames per clip $T$             & $24$ \\
Backbone                        & ResNet50, ImageNet pre-trained, fine-tuned \\
SE reduction ratio              & $16$ \\
BiLSTM layers                   & $2$ \\
BiLSTM hidden size              & $256$ per direction \\
Temporal attention hidden size  & $128$ \\
Classifier hidden size          & $256$ \\
Dropout                         & $0.5$ \\
Loss                            & BCE with logits, per-label \texttt{pos\_weight} clipped to $[0.5, 6.0]$ \\
Optimiser                       & AdamW (\texttt{lr}\,$=5\!\times\!10^{-5}$, \texttt{weight\_decay}\,$=10^{-5}$) \\
Gradient clipping               & global norm $\leq 5.0$ \\
LR scheduler                    & \texttt{ReduceLROnPlateau} on val.\ Macro F1 (factor $0.5$, patience $2$) \\
Batch size                      & $4$ \\
Epochs                          & up to $12$, early stop patience $4$ on val.\ Macro F1 \\
Threshold grid                  & $0.10, 0.15, \dots, 0.70$ per label, per-label F1 max \\
Train-time augmentation         & flip $0.5$, brightness/contrast $0.5$, rotation $0.25$, blur $0.15$, crop $0.25$ \\
Oversampling multipliers        & engagement $\times 1$, boredom $\times 2$, confusion $\times 3$, frustration $\times 4$ \\
\bottomrule
\end{tabular}
\end{table}

\section{Integration Readiness for Adaptive Learning}\label{sec:method_integration}

The FER module is designed to be a clean input source for the adaptation module, not an end product on its own. At inference time, the module takes a $24$-frame webcam clip and returns a four-dimensional emotion vector, where each component is the per-label sigmoid score for boredom, engagement, confusion, and frustration. These scores are continuous values between $0$ and $1$, computed with the same per-label thresholds as during evaluation, which makes the output stable across clips. The adaptation module described in the next section consumes this vector together with simple performance signals, such as whether the previous answer was correct, to choose the next teaching action and the next question.

This separation between perception and decision has practical advantages. It allows the FER module to be retrained or replaced without changing the adaptation logic, as long as the output interface remains a four-dimensional emotion vector. It also allows the adaptation module to be evaluated under different emotional input conditions, including the case where the FER module is replaced by a no-emotion baseline, which is the comparison setup used later in this thesis. Finally, because the FER module only requires a webcam, the integrated system inherits the same lightweight deployment profile, which is in line with the FER-only adaptive design promoted by recent educational FER reviews \cite{lek_academic_2023,florestiyanto_emotion_2024,vistorte_integrating_2024}. The exact way this emotion vector is fed into the adaptation policy, and how the policy is trained and evaluated, is described in the following section.
